\cvsectionnew{Expériences Professionnelles}
\cveventnew{Stage de Fin d'Etudes}{PwC Maroc }{Février 2024 -- Août 2024 \hspace{1 cm}}{}
\vspace{-4mm}
\textbf{Développement d'une infrastructure Red Teaming}
\begin{itemize}
    \setlength\itemsep{0cm}
    \item Étude des exigences de l'équipe Red Team et création de l'architecture.
    \item Réalisation de l'architecture sur le cloud privé de PwC et AWS.
    \item Intégration du logging et du monitoring sur l'infrastructure.
    \item Tests sur une machine victime.
\end{itemize}
\vspace{-1mm}
\textit{Technologies: Sliver, Gophish, ELK Stack, Terraform, AWS, Socat, Docker}
% \bigskip
\vspace{3mm}

\cveventnew{Stage Technique}{INPT }{ Juin 2023 -- Septembre 2023 \hspace{1 cm}}{}
\vspace{-4mm}
\textbf{Analyse des techniques pour contourner les protections de "packing" dans les malwares}
\begin{itemize}
    \setlength\itemsep{0cm}
    \item Etude sur la structure de l'exécutable Windows et création d'un packer.
    \item Découverte des techniques de "unpacking".
    \item Étude de cas sur deux malwares. (Brbbot, Trickbot)
\end{itemize}
\vspace{-1mm}
\textit{Technologies: pestudio, procmon, tcpmon, Regshot, API Monitor, Ghidra, DiE, CFF Explorer, x64dbg, PE-Bear, pestudio, Cutter, windbg, C++, API Windows, zlib}